\chapter{Pen-and-Paper Security Proof Reconstructed from EasyCrypt}
\label{ch:pen-paper-reconstruction}

This chapter reconstructs the pen-and-paper proof of \haetae{} provable security
using the EasyCrypt development as a checked source of lemmas.  The aim is not
to analyze the proof script line by line.  The aim is to write the mathematical
hybrid argument that a cryptographer would expect on paper, while citing the
EasyCrypt theorem names that certify each step.

\section{Proof convention}
\label{sec:paper-proof-convention}

Let \(A\) be a probabilistic polynomial-time adversary in the modeled
random-oracle interface.  The adversary interacts with a hash oracle \(H\) and a
signing oracle \(O\), then outputs a candidate forgery
\((m^\star,ctx^\star,\sigma^\star)\).  The event
\(\mathsf{Forge}\) is the event that verification accepts and the message-context
pair was not previously signed.

For the informal target cryptographic game, we write
\[
  \Adv_{\haetae,\rom}^{\eufcma}(A)
  =
  \Pr[\mathsf{Forge}\text{ in the real EUF-CMA game}].
\]
The checked theorem reconstructed below uses the EasyCrypt-model advantage; the
fixed query constants belong to the proof ledger and bound terms rather than to a
syntactic restriction on the source \eufcma{} game.
All probabilities in this chapter are over the coins of key generation, signing,
the random oracle, and the adversary.  When this reconstruction mirrors
paper-style notation, \(q_H\) and \(q_S\) denote the fixed wrapper counts used by
the proof ledger; in the EasyCrypt development these are represented by
\ec{hash\_query\_budget\_count} and \ec{signature\_query\_budget\_count}.

The concrete counted-ROM loss expression is the one defined in
\ec{HAETAE\_Reductions.ec}.  For the theorem emphasized in this dissertation we
write it schematically as
\[
  B_{\haetae}
  =
  B_{\mathsf{hard}} + B_{\mathsf{ROM}},
\]
where \(B_{\mathsf{hard}}\) expands to the MLWE, Module-SIS, and
bimodal-to-MSIS terms, and \(B_{\mathsf{ROM}}\) expands to the counted collision,
prequery, and min-entropy terms.  The paper-style
\ec{haetae\_euf\_bound} contains a separate rejection/Fiat--Shamir-with-aborts
ledger, but it is not the final counted-ROM theorem.  The counted-ROM grouping is
certified by
\ec{haetae\_euf\_counted\_rom\_bound\_groupedE}.

\begin{theorem}[Reconstructed formal-model counted-ROM implication]
\label{thm:reconstructed-haetae-euf}
For every adversary \(A\) satisfying the declared EasyCrypt module interfaces,
with the proof and bound ledger instantiated through the fixed wrapper constants
\(q_H=q_S=2\), if the explicit hop/reduction premise and the MLWE and
Module-SIS hardness premises required by
\ec{euf\_cma\_security\_from\_checked\_ro\_sampling\_hop\_and\_counted\_bimodal}
hold for the declared reduction modules appearing in those antecedents, then the
formal EasyCrypt model satisfies
\[
  \Adv_{\mathsf{ECModel},\rom}^{\eufcma}(A)
  \le
  \ec{haetae\_euf\_counted\_rom\_bound}.
\]
Equivalently, after expanding the counted-ROM bound,
\[
  \Adv_{\mathsf{ECModel},\rom}^{\eufcma}(A)
  \le
  L_{\mathsf{MLWE}}
  +
  L_{\mathsf{ModSIS}}
  +
  L_{\mathsf{bimodal}\to\mathsf{MSIS}}
  +
  L_{\mathsf{ROM}},
\]
where the hardness terms and ROM loss terms are the concrete real expressions in
\ec{HAETAE\_Reductions.ec}.  This is not a specification-level theorem for
\(\mathsf{Spec}_{260204}\), and it is not an arbitrary-query theorem.
\end{theorem}

\begin{proof}
The proof is a sequence of games.  Each displayed transition is an ordinary
paper-proof step; the cited EasyCrypt lemmas are the machine-checked
certificates that the step is valid.  The final composition is not a bare
closed theorem over all adversaries: it is a fixed-ledger formal-model
implication whose antecedents are the final NMA/hardness hop and the declared
MLWE and Module-SIS bounds.
\end{proof}

\section{Correctness and the starting game}
\label{sec:paper-proof-starting-game}

The starting game \(G_0\) is the standard, unbudgeted EUF-CMA game for the
random-oracle signature interface:
\[
  G_0 = \ec{Sig\_ROM.EUF\_CMA}(\haetae,H,A).
\]
Thus, inside the EasyCrypt source game used by the checked composition,
\[
  \Pr[G_0(A)\Rightarrow 1]
  =
  \Adv_{\mathsf{ECModel},\rom}^{\eufcma}(A).
\]
The fixed constants \(q_H=q_S=2\) enter through the later budgeted wrappers and
the right-hand-side bound, not through this source-game notation.
The correctness side of the theorem is independent of the unforgeability
hybrid.  It states that honestly generated signatures verify:
\[
  \Pr[
    \mathsf{Verify}(pk,m,ctx,\mathsf{Sign}(sk,m,ctx))=1
  ]=1.
\]
The EasyCrypt proof anchors are
\ec{correctness\_obligation\_holds} and
\ec{correctness\_perfect}.  They discharge the paper proof's statement that the
modeled signing and verification procedures are consistent.

\section{Removing the direct signing view}
\label{sec:paper-proof-simulated-signing}

The first unforgeability step replaces real signing calls by a signing simulator
that exposes the same distributional interface to the adversary.  Let \(G_1\)
be the simulated-signing game:
\[
  G_1 = \ec{EUF\_CMA\_SimulatedSign}.
\]
The paper claim is that the adversary's success probability is unchanged, or
bounded by the explicit signing-hop loss when one moves through the more refined
interfaces:
\[
  \Pr[G_0(A)\Rightarrow 1]
  \le
  \Pr[G_1(A)\Rightarrow 1] + L_{\mathsf{sign}}.
\]
The corresponding checked statements are
\ec{euf\_cma\_security\_from\_simulated\_hop},
\ec{euf\_cma\_security\_from\_explicit\_boundaries}, and
\ec{euf\_cma\_security\_with\_hop\_interfaces}.

At the paper level, this step is justified by the fact that a signature is fully
determined by the transcript fields sampled by the signing algorithm.  The
transcript predicates record the required public equation, challenge query,
response, hint, and rejection-branch consistency.  These facts are certified by
\ec{transcript\_valid\_public\_equation},
\ec{transcript\_valid\_signature\_challenge\_query},
\ec{transcript\_fields\_match\_signature\_challenge\_query}, and
\ec{transcript\_valid\_commitment\_reconstruction}.

\section{Internal transcript and random-oracle view}
\label{sec:paper-proof-rom-internal}

The next step internalizes the transcript log and the random-oracle state.  Let
\(G_2\) be the transcript-simulated signing game, \(G_3\) the internal
transcript game, and \(G_4\) the ROM-internal transcript game:
\[
  G_2 = \ec{EUF\_CMA\_TranscriptSimulatedSign},
  \qquad
  G_3 = \ec{EUF\_CMA\_InternalTranscriptSign},
  \qquad
  G_4 = \ec{EUF\_CMA\_ROMInternalTranscriptSign}.
\]
The paper-proof statement is
\[
  \Pr[G_1(A)\Rightarrow 1]
  \le
  \Pr[G_4(A)\Rightarrow 1] + L_{\mathsf{tr}}.
\]
The reason is that transcript logging and lazy-ROM internalization do not change
the adversary's observable interface, except at explicitly counted bad events.

The EasyCrypt invariants used here are:
\[
\begin{array}{ll}
\ec{transcript\_log\_consistent} &
  \text{logged signatures match transcript fields},\\
\ec{transcript\_log\_ready} &
  \text{transcripts are ready for programming},\\
\ec{internal\_transcript\_state\_sound} &
  \text{internal state agrees with key and transcript data},\\
\ec{sampler\_rom\_covered} &
  \text{sampler ROM calls are represented in the ROM state}.
\end{array}
\]
Their preservation under signing is checked by
\ec{transcript\_log\_consistent\_sign\_internal\_cons},
\ec{internal\_transcript\_state\_sound\_sign\_internal\_cons}, and the
sampler-coverage lemmas
\ec{sampler\_rom\_covered\_fresh\_after\_clean\_seed},
\ec{sampler\_rom\_covered\_self\_log\_preserves}, and
\ec{sampler\_rom\_covered\_old\_log\_clean\_boundary}.

In paper terms, this is the point where the proof stops treating the random
oracle as an external black box and begins treating it as a lazy table whose
queries, programmed sites, and bad events can be counted.

\section{ROM programming and counted bad events}
\label{sec:paper-proof-rom-programming}

The Fiat--Shamir-with-aborts proof programs challenge sites in the random
oracle.  A programming site is a pair
\[
  s=(q,y)\in \mathsf{Q}_{\mathsf{RO}}\times\mathsf{Y}_{\mathsf{RO}},
\]
where \(q\) is the query to be programmed and \(y\) is the desired output.  The
programming is sound when the site is fresh:
\[
  q\notin Q_A,
\]
where \(Q_A\) is the adversary's previous random-oracle query log.

The paper proof decomposes the failure probability into counted events:
\[
  L_{\mathsf{ROM}}
  =
  L_{\mathsf{coll}}
  +
  L_{\mathsf{pre}}
  +
  L_{\mathsf{me}}.
\]
This corresponds to the EasyCrypt terms
\ec{counted\_rom\_collision\_term},
\ec{counted\_rom\_prequery\_term}, and
\ec{counted\_rom\_min\_entropy\_term}.  Their grouping is certified by
\ec{counted\_rom\_programming\_loss\_termE} and
\ec{counted\_rom\_programming\_loss\_splitE}.

The preservation laws for the lazy ROM are the checked form of the paper
argument that lazy sampling remains consistent after query, observation, and
programming:
\[
\begin{array}{ll}
\ec{counted\_lazy\_rom\_init\_clears} &
  \text{initial state has no bad event},\\
\ec{counted\_lazy\_rom\_get\_preserves\_clear} &
  \text{ordinary queries preserve cleanness},\\
\ec{counted\_lazy\_rom\_program\_preserves\_bad\_clear\_if\_fresh} &
  \text{fresh programming does not create badness},\\
\ec{counted\_lazy\_rom\_bad\_false\_when\_clear} &
  \text{the clear invariant excludes the bad event}.
\end{array}
\]

\section{From chosen-message to no-message structure}
\label{sec:paper-proof-nma}

After transcript and ROM internalization, the signing oracle is represented by
logged samples.  The proof then factors the active signing interaction into a
no-message attack against a simulated transcript distribution.  At the paper
level this says that the signing oracle has been replaced by a sampler whose
outputs are independent of the final forgery except through the public
transcript and ROM state.

Let \(G_5\) be the ROM-internal NMA game obtained from the chosen-message game:
\[
  G_5 = \ec{ROMInternalTranscriptAsNMA}.
\]
The generic exact adapter is certified by \ec{nma\_as\_cma\_exact}, and the
\haetae{}-specific chain is certified by
\ec{euf\_cma\_security\_from\_structural\_nma\_and\_counted\_bimodal},
\ec{euf\_cma\_security\_from\_public\_sim\_nma\_and\_counted\_bimodal}, and
\ec{euf\_cma\_security\_from\_paper\_sim\_nma\_and\_counted\_bimodal}.

Thus the paper inequality is
\[
  \Pr[G_4(A)\Rightarrow 1]
  \le
  \Pr[G_5(A')\Rightarrow 1] + L_{\mathsf{NMA}},
\]
where \(A'\) is the induced no-message adversary.

\section{Rejection sampling and the paper sampler}
\label{sec:paper-proof-rejection-sampling}

The paper proof next replaces the real rejection-sampling view by the paper
sampler view.  This step is justified by the fact that the signing attempt state
contains all fields needed to reconstruct an accepted signature:
\[
  (y_1,y_2,w,c,z,h,\rho)
  \longmapsto
  \sigma.
\]
The paper claim is that the sampler distributions are related up to the
rejection-sampling loss:
\[
  \Pr[G_5(A')\Rightarrow 1]
  \le
  \Pr[G_6(A')\Rightarrow 1] + L_{\mathsf{rej}}.
\]
The checked anchors are
\ec{signing\_attempt\_verifies},
\ec{signing\_simulation\_verifies},
\ec{signing\_attempt\_simulates},
\ec{structural\_signing\_sample\_pair\_sampler\_ok},
\ec{bounded\_unit\_signing\_sample\_pair\_sampler\_ok},
\ec{checked\_hyperball\_signing\_sample\_pair\_sampler\_ok}, and
\ec{exact\_hyperball\_signing\_sample\_pair\_sampler\_ok}.

In the top-level chain this is reflected by
\ec{euf\_cma\_security\_from\_real\_signing\_paper\_sim\_nma\_and\_counted\_bimodal},
\ec{euf\_cma\_security\_from\_haetae\_rejection\_paper\_sim\_nma\_and\_counted\_bimodal},
and
\ec{exact\_hyperball\_haetae\_rejection\_nma\_bound\_from\_paper\_sim}.

\section{Exact hyperball sampling}
\label{sec:paper-proof-exact-hyperball}

The exact hyperball game replaces the checked paper sampling distribution by the
ideal exact distribution on the hyperball.  The mathematical step is
\[
  \Pr[G_6(A')\Rightarrow 1]
  \le
  \Pr[G_7(A')\Rightarrow 1] + L_{\mathsf{hyp}}.
\]
The distributional facts are encoded in
\ec{HAETAE\_Distributions.ec}; the proof uses losslessness, support, and point
bound lemmas for the relevant signing sample distributions.

The security chain exposes this step through
\ec{haetae\_rejection\_nma\_bound\_from\_exact\_hyperball\_sampling\_hop},
\ec{haetae\_rejection\_exact\_hyperball\_sampling\_hop\_from\_real\_signing\_exact\_hyperball\_hop},
and
\ec{haetae\_rejection\_ro\_exact\_hyperball\_sampling\_hop\_from\_real\_signing\_ro\_exact\_hyperball\_hop}.

\section{Tight RO-signing-attempt to RO-exact-hyperball lifting}
\label{sec:paper-proof-budgeted-lifting}

The most delicate part of the proof is the transition from the RO signing
attempt sampler to the RO exact hyperball sampler.  Let \(R\) be fresh signing
randomness sampled from \(\mu_{\mathsf{coins}}\), and define the sampler query
\[
  q_R = \ec{sampler\_expand\_query}(R).
\]
Let \(Q_A\) be the adversary's logged hash-query set.  The clean event is
\[
  \mathsf{Clean}
  =
  \{q_R\notin Q_A\}
  \wedge
  \{q_R\notin Q_{\mathsf{sampler}}\}
  \wedge
  \mathsf{sampler\_rom\_covered}
  \wedge
  \mathsf{budget\ discipline}.
\]
Here \(Q_{\mathsf{sampler}}\) is the self-log of sampler-expansion queries from
earlier signing calls.  The checked event \ec{sampler\_bad\_prequery} is set if
the fresh sampler query is already in either log.
On \(\mathsf{Clean}\), one signing call in the two games is coupled so that its
visible output is identical.  Therefore the local fundamental-lemma contribution
of that call is bounded by
\[
  \Pr[q_R\in Q_A\cup Q_{\mathsf{sampler}}].
\]
Since \(R\) is fresh, \(Q_A\) contains at most \(q_H\) logged adversary hash
queries, and \(Q_{\mathsf{sampler}}\) contains at most \(q_S\) sampler self-log
queries over the whole budgeted execution, the one-call loss is
\[
  \Pr[q_R\in Q_A\cup Q_{\mathsf{sampler}}]
  \le
  (q_H+q_S)\cdot
  \max_r \Pr[R=r].
\]
The displayed inequality is a local signing-call bound.  The full \(G_7\) to
\(G_8\) hop must also union over the at most \(q_S\) signing calls, so the
full-game sampler-prequery contribution is
\[
  \left|
  \Pr[G_7(A')\Rightarrow 1]
  -
  \Pr[G_8(A')\Rightarrow 1]
  \right|
  \le
  q_S(q_H+q_S)\cdot
  \max_r \Pr[R=r].
\]
In the EasyCrypt development the point bound on \(R\) is supplied by the
signing-randomness distribution lemmas and the budgeted support lower bound,
yielding the checked full-game term
\[
  \ec{rom\_signature\_query\_budget}
  \cdot
  \frac{\ec{rom\_hash\_query\_budget}+
        \ec{rom\_signature\_query\_budget}}
       {\ec{challenge\_support\_cardinality\_lower\_bound}}.
\]
The NMA lifting lemmas separately carry
\(\ec{rom\_signature\_query\_budget}\cdot
\ec{rejection\_sampling\_loss\_term}\); this is distinct from the
sampler-prequery loss above.

The checked anchors are
\ec{rom\_internal\_budgeted\_nma\_ro\_signing\_attempt\_ro\_exact\_hyperball\_clean\_conditioned\_lifting},
\ec{rom\_internal\_budgeted\_nma\_ro\_signing\_attempt\_ro\_exact\_hyperball\_loss\_from\_paper\_sample\_lifting},
\ec{concrete\_budgeted\_self\_logged\_signature\_clean\_event\_loss\_surface},
\ec{concrete\_budgeted\_o\_sign\_clean\_one\_call\_loss\_from\_self\_logged\_surface},
and the active projection lemmas for the attempt side and exact side.

The top-level security file uses this lifting through
\ec{real\_signing\_ro\_exact\_hyperball\_hop\_from\_budgeted\_lifting},
\ec{haetae\_rejection\_ro\_exact\_hyperball\_sampling\_hop\_from\_budgeted\_lifting},
\ec{real\_signing\_ro\_exact\_hyperball\_hop\_from\_paper\_sample\_lifting}, and
\ec{euf\_cma\_security\_from\_budgeted\_paper\_sample\_lifting\_and\_counted\_bimodal}.

\section{Forgery extraction and hardness reduction}
\label{sec:paper-proof-hardness}

In the final exact-hyperball NMA game, extraction is a forking-pair
special-soundness step, not a conclusion from one valid transcript alone.  When
the preceding game and reduction interfaces provide two compatible accepting
transcripts with the required shared public data and distinct challenges, the
local extraction obligation has the following form:
\[
  \mathsf{ValidForkingPair}(\mathsf{tr}_1,\mathsf{tr}_2)
  \Longrightarrow
  \mathsf{Solve}_{\mathsf{Bimodal/MSIS}}.
\]
The transcript-to-relation facts are checked by
\ec{extract\_bimodal\_forking\_some},
\ec{extract\_bimodal\_forking\_nonempty},
\ec{valid\_forking\_pair\_public\_equations}, and
\ec{bimodal\_to\_module\_sis\_valid}.

The reduction adapters are exact in the sense that they preserve the success
event of the induced solver:
\[
  \Pr[\mathsf{BimodalSolver}(A)\Rightarrow 1]
  \le
  \Pr[\mathsf{ModuleSIS}(B_A)\Rightarrow 1].
\]
The checked anchors are
\ec{bimodal\_solver\_lift\_exact},
\ec{bimodal\_mode\_solver\_lifted\_distribution\_exact}, and
\ec{bimodal\_msis\_as\_module\_sis\_exact}.  The top-level composition uses
\ec{bimodal\_to\_module\_sis\_reduction\_bound} and the counted-bimodal security
lemmas in \ec{HAETAE\_Security.ec}.

\section{Composing the inequalities}
\label{sec:paper-proof-composition}

Combining the checked counted-ROM path under the final hop/reduction premise and
the MLWE and Module-SIS hardness premises gives
\[
\begin{aligned}
  \Adv_{\mathsf{ECModel},\rom}^{\eufcma}(A)
  &=
  \Pr[G_0(A)\Rightarrow 1] \\
  &\le
  \Pr[G_8(A')\Rightarrow 1] + L_{\mathsf{ROM}} \\
  &\le
  L_{\mathsf{MLWE}}
  + L_{\mathsf{ModSIS}}
  + L_{\mathsf{bimodal}\to\mathsf{MSIS}}
  + L_{\mathsf{ROM}}
\end{aligned}
\]
Here \(L_{\mathsf{ROM}}\) is the counted collision, prequery, and min-entropy
loss.  Rejection, exact-sampling, and active-signing lifting obligations are real
proof obligations in the hop chain, but in the counted theorem they are pushed
into explicit hypotheses and bridge lemmas rather than appearing as additional
final summands.  This is not the same as proving that the hypotheses are
non-vacuously satisfiable.  In the current EasyCrypt constants,
\ec{rejection\_sampling\_loss\_term} is already larger than one, so a final
antecedent containing
\(\Pr[\cdot]+\ec{rejection\_sampling\_loss\_term}\le \Pr[\cdot]+\Pr[\cdot]\)
must be treated as a vacuity warning until the rejection ledger is normalized or
the antecedent is replaced by a stronger checked hop.  The separate expression
\ec{haetae\_euf\_bound} records the older paper-style
rejection/Fiat--Shamir-with-aborts ledger; it should not be added to
\ec{haetae\_euf\_counted\_rom\_bound}.
The real-arithmetic side conditions required to make this a probability-scale
bound are checked in \ec{HAETAE\_Reductions.ec}:
\[
\begin{array}{ll}
\ec{counted\_rom\_programming\_loss\_term\_nonnegative} &
  L_{\mathsf{ROM}}\ge 0,\\
\ec{counted\_rom\_programming\_loss\_term\_subunit} &
  L_{\mathsf{ROM}}\le 1,\\
\ec{challenge\_support\_cardinality\_lower\_bound\_positive} &
  |\mathcal{C}|>0,\\
\ec{counted\_rom\_support\_dominates\_budget} &
  \text{the support lower bound dominates the counted query budget},\\
\ec{counted\_rom\_loss\_from\_query\_budgets\_and\_support} &
  \text{the counted ROM loss follows from the query and support bounds}.
\end{array}
\]

Finally, \ec{euf\_cma\_security\_from\_checked\_ro\_sampling\_hop\_and\_counted\_bimodal}
and \ec{euf\_cma\_security\_from\_real\_signing\_exact\_hyperball\_hop\_and\_counted\_bimodal}
compose the checked sampling hop with the counted bimodal reduction.  This is
the fixed-ledger formal-model implication reconstructed at the beginning of the
chapter, subject to the explicit antecedents and vacuity warning above.

\section{What this reconstruction does and does not claim}
\label{sec:paper-proof-scope}

This chapter is a pen-and-paper reconstruction of the proof.  It does not add a
new EasyCrypt theorem.  The machine-checked content remains the EasyCrypt files
listed in the provable-security manifest.  The contribution of this chapter is
expository: it recovers the mathematical hybrid proof from the checked lemma
graph and records which EasyCrypt lemma certifies each paper step.

The proof remains conditional on the intended interpretation of the abstract
hardness games and on the fidelity of the HAETAE model to the specification.
Within those boundaries, the game hops, sampler transformations, ROM
programming steps, bad-event bounds, and final arithmetic composition are
machine checked by the cited EasyCrypt development.
